% GENERATED FROM wb-ref-wall.md BY md2tex.py -- DO NOT HAND-EDIT.
\documentclass[11pt,a4paper]{article}
\usepackage[margin=1in]{geometry}
\usepackage{amsmath,amssymb,bm}
\usepackage{listings}
\usepackage[hidelinks]{hyperref}
\usepackage{longtable}
\usepackage{booktabs}
\setlength{\parskip}{0.6em}
\setlength{\parindent}{0pt}
\lstset{basicstyle=\ttfamily\footnotesize,breaklines=true,frame=single,
        columns=fullflexible,keepspaces=true}
\title{The default well-balanced x1 reference at a physical wall}
\author{snapy --- derivation}\date{}
\begin{document}
\maketitle
\section*{The default well-balanced x1 reference at a physical wall}

This is the default reference, with \texttt{SNAP\_WB\_REF4} off and \texttt{dynamics/wb-wall-clamp} on (the default). With the
repeated wall cell, the face density reference has an $O(\Delta z)$ error at the first three faces next to each
wall, against $O(\Delta z^2)$ elsewhere. Continuing $\rho/p$ past the wall, linearly or ln-linearly, restores the
interior order there and leaves every other face bit for bit. The code is \texttt{src/hydro/hydro\_ref\_x1\_impl.h}
(\texttt{hydro\_ref\_x1\_rop\_smooth}, \texttt{hydro\_ref\_x1\_cell\_impl}), with the tensor (MPS) path in
\texttt{src/hydro/hydro\_dispatch.cpp}. The test is \texttt{tests/test\_wb\_ref\_wall.cpp}.

\section*{1. What the kernel builds}

Per column, with $r_i = \bar\rho_i/\bar p_i$ and the binomial $B = (1, 4, 6, 4, 1)/16$:

\begin{equation*}
r^s_i = \sum_{m=-2}^{2} B_m\, r_{i+m},\qquad
\rho_{{\rm ref},i} = p_{{\rm ref},i}\, r^s_i,\qquad
\rho_{{\rm sf},f} = p_{{\rm sf},f}\,\tfrac12\big(r^s_{f-1} + r^s_f\big),
\end{equation*}

where face $f$ is the lower face of cell $f$. The solver's face density is $\rho_{\rm sf} + \mathcal W[\bar\rho -$
\textbackslash\{\}rho\_\{\textbackslash\{\}rm ref\}]$, where $\textbackslash\{\}mathcal W$ is WENO5 and the perturbation's ghosts at a wall are filled with even parity$
(\texttt{hydro\_forward.cpp}). At 37dce4e:
\begin{itemize}
\item the clamp is at \texttt{hydro\_ref\_x1\_impl.h:74}: an index past the wall is replaced by the wall cell;
\item \texttt{jlo}/\texttt{jhi} are set at \texttt{:159-160};
\item $r^s$, $\rho_{\rm ref}$ and $\rho_{\rm sf}$ are formed at \texttt{:161-167};
\item the even-parity fill is at \texttt{hydro\_forward.cpp:276-292}.
\end{itemize}

\section*{2. Interior order}

For smooth $r$, $B$ has unit sum, zero first moment and second moment 1 (in cells), so $r^s_i = r_i +$
\textbackslash\{\}tfrac\{\textbackslash\{\}Delta z\textasciicircum{}2\}\{2\} r''\_i + O(\textbackslash\{\}Delta z\textasciicircum{}4)$. The mean of two neighbouring cells adds $\textbackslash\{\}tfrac\{\textbackslash\{\}Delta z\textasciicircum{}2\}\{8\}r''$.$
With $p_{{\rm sf},f} = p(z_f)$ from the scan, and $r_i = R_i + O(\Delta z^2)$ for a ratio of averages,

\begin{equation*}
\rho_{{\rm sf},f} = \rho(z_f) + O(\Delta z^2)\qquad\text{(interior faces)}.
\end{equation*}

\section*{3. The repeated wall cell is first order}

Near the bottom wall, write $r_j = r_0 + j\,a + O(\Delta z^2)$ with $a = r'\Delta z = O(\Delta z)$, and index the
cells from the wall, $j = 0, 1, \dots$. Repeating $r_0$ for every $j < 0$ gives:

\begin{itemize}
\item wall cell: $r^s_0 = (11 r_0 + 4 r_1 + r_2)/16 = r_0 + \tfrac{3}{8}a$;
\item next cell: $r^s_1 = (5 r_0 + 6 r_1 + 4 r_2 + r_3)/16 = r_1 + \tfrac{1}{16}a$;
\item first ghost (it enters the wall face): $r^s_{-1} = r_0 + \tfrac{1}{16}a$, where the true value is $r_0 - a$.
\end{itemize}

Against the true face values $r(z_f)$, the face references are off by:

\begin{longtable}{ll}
\toprule
face (cells from the wall) & error of $\tfrac12(r^s_{f-1}+r^s_f)$ \\
\midrule
0 (the wall face) & $\tfrac{23}{32}a$ \\
1 & $\tfrac{7}{32}a$ \\
2 & $\tfrac{1}{32}a$ \\
$\ge 3$ & $O(\Delta z^2)$ \\
\bottomrule
\end{longtable}

At the top wall the signs flip. Relative to $r$ the error is a fixed multiple of $(r'/r)\Delta z$, so it is
\textbf{first order}. It vanishes when $r' = 0$, so an isothermal column has no wall error.

Check against the test column (polytrope $T = 1 - \beta z$, $\beta = 0.5$, one pressure scale height, nz 64,
$\Delta z = 0.01230$; $r'/r = \beta/T$). The prediction at face 1 is $\tfrac{7}{32}\cdot0.5\cdot0.0123 = 1.35\times10^{-3}$,
and \texttt{test\_wb\_ref\_wall} measures $1.39\times10^{-3}$. At face 2 the prediction is $1.9\times10^{-4}$ plus the
interior $O(\Delta z^2)$ part; it measures $2.5\times10^{-4}$.

\textbf{The even-parity ghosts.} $\rho' = \bar\rho - \rho_{\rm ref}$ inherits the wall bias,
$-\tfrac38 a\,p$ in the wall cell and $-\tfrac1{16}a\,p$ in the next. The even mirror carries that $O(\Delta z)$
feature into the ghosts, so $\mathcal W[\rho']$ does not cancel it, and the solver's face densities $\rho_{L,R}$ at
faces 1 and 2 are also first order (measured: $\rho_L$ at face 1 is $-5.2\times10^{-4}$, then $-2.5\times10^{-4}$ at nz 32, 64).
The mirror is not the source. Applied to a $\rho'$ that is $O(\Delta z^2)$ and smooth near the wall, it costs
only $O(\Delta z^2)$. The source is the bias in $\rho_{\rm ref}$.

\section*{4. The closure: continue $\rho/p$ past the wall, linearly or ln-linearly}

Past a clamped wall, take

\begin{equation*}
r_{-k} = \begin{cases} r_0 + k\,(r_0 - r_1), & r_1 \le r_0,\\ r_0\,(r_0/r_1)^k, & r_1 > r_0,\end{cases}
\qquad k = 1, 2, 3
\end{equation*}

(and the mirror image at the top), in place of $r_0$ (\texttt{hydro\_ref\_x1\_wall\_rop}). The first branch is the linear
continuation of $r$ and the second that of $\ln r$. The two differ by $O(\Delta z^2)$, and $B$ reproduces a linear
profile exactly, so the $O(a)$ terms above cancel. What remains is the continuation's own curvature error,
$\tfrac{k(k+1)}{2}\Delta z^2$ times $r''$ or $(\ln r)''\,r$, which is $O(\Delta z^2)$. The wall cells, and faces
0–2, then have the interior's order. For constant $r$ the first branch gives $r_0$ exactly, so an isothermal
column is unchanged bit for bit.

\textbf{Why two branches.} Each branch is unbounded on one side:
\begin{itemize}
\item linear in $r$ goes to zero when $r$ rises away from the wall. With $d = (r_1 - r_0)/r_0$, $r_{-3} = r_0(1 - 3d)$
  vanishes at $d = 1/3$ and is negative past it;
\item ln-linear goes to infinity when $r$ falls away from the wall. A nearly empty neighbour, $r_1 \ll r_0$, gives
  $r_0(r_0/r_1)^k$.
\end{itemize}

Each case takes the branch that stays closer to $r_0$, so for $r_1 \ge 0$ the continuation lies in
$[\,r_0(r_0/r_1)^k,\; r_0(1+k)\,]$. In \texttt{test\_face\_floor}'s unresolved column, the cell next to the top wall is
dipped to $10^{-4}$, and the dipped-face mass flux is
\begin{itemize}
\item $3.0\times10^{-8}$ with the clamp,
\item $2.83\times10^{-8}$ with either the linear form or this closure,
\item $3.36\times10^{-7}$ with the ln-linear form everywhere.
\end{itemize}

Guards (the clamp stays):
\begin{itemize}
\item the wall side must own at least two cells, so only owned cells are read and the clamp's property tests in
  \texttt{test\_hydro\_ref\_x1} hold;
\item a value that is not positive falls back to $r_0$.
\end{itemize}

Measured with \texttt{test\_wb\_ref\_wall}, $\beta = 0.5$, dsf at face 1 (the bottom wall, where $r$ rises away from the
wall, so the ln-linear branch applies; the top wall takes the linear branch):

\begin{longtable}{llll}
\toprule
nz & clamp only & linear in $r$ everywhere & this closure \\
\midrule
16 & $6.16\times10^{-3}$ & $6.48\times10^{-4}$ & $8.04\times10^{-4}$ \\
32 & $2.88\times10^{-3}$ & $1.57\times10^{-4}$ & $1.95\times10^{-4}$ \\
64 & $1.39\times10^{-3}$ & $3.87\times10^{-5}$ & $4.82\times10^{-5}$ \\
128 & $6.84\times10^{-4}$ & $9.78\times10^{-6}$ & $1.22\times10^{-5}$ \\
\bottomrule
\end{longtable}

The observed order goes from about 1.05 to about 2.0 (1.81–2.12 over faces 1–2 at both walls, for dsf,
$\rho_L$ and $\rho_R$, nz 32 → 64). The solver's $\rho_L$ at face 1 goes from $-2.54\times10^{-4}$ (nz 64) to
$7.5\times10^{-6}$. Faces 1–2 at both walls stay at or below the largest interior error ($1.45\times10^{-4}$ at
nz 64).

\textbf{What changes.} At each clamped physical wall: $\rho_{\rm sf}$ at faces 0, 1, 2 and $\rho_{\rm ref}$ at cells 0, 1,
plus their ghost rows. $p_{\rm sf}$, $p_{\rm ref}$ and every other face and cell are bit for bit as before
(compared at nz 16–128). The rest balance does not involve $\rho_{\rm sf}$ or $\rho_{\rm ref}$ (§1 of
\texttt{wb-ref4.md}), so \texttt{balance\_column} and the discrete rest state are unchanged.

\textbf{Why not the \texttt{SNAP\_WB\_REF4} closure.}
\begin{itemize}
\item \texttt{SNAP\_WB\_REF4}'s wall closure (\texttt{wb\_ref4.cpp}) is a post-pass that replaces $\rho_{\rm ref}$ with
  $p_{\rm ref}F(r)$ and $\rho_{\rm sf}$ with a fourth-order face value on *every* face, so applying it would change
  every interior face. Its cubic continuation is tied to $F$: $F$ with the cubic values returns $r$ at the end
  cells exactly. It also needs four owned cells.
\item A cubic (or quadratic) continuation inside $B$ would also restore second order, with weights of 10–20 against
  $1+k$; it was not chosen.
\end{itemize}

\section*{5. The straka CFL 1.6 run}

\texttt{test\_straka\_redo} runs \texttt{examples/straka\_single.yaml} with CFL 1.6 and tlim 60 s, and fails if one step needs more
than 5 redos. CFL 1.6 is past the scheme's stability limit, so this is a robustness test, and it is marginal: most
redos are triggered by non-finite states. Runs reaching 60 s, with nx1 64, one thread, and
$\Delta T = -15(1 + k\epsilon)$:

\begin{longtable}{llll}
\toprule
reference & test deck & $\epsilon = 10^{-12}$, $k = 0..9$ & $\epsilon = 10^{-4}$, $k = 0..19$ \\
\midrule
clamp only & pass & 10/10 & 20/20 \\
linear in $r$ everywhere & fail (31.67 s) & 0/10 & 18/20 \\
ln-linear everywhere & pass & 9/10 & 19/20 \\
this closure & pass & 10/10 & 18/20 \\
\bottomrule
\end{longtable}

With $\epsilon = 10^{-12}$ the runs of this closure stay on one trajectory, ending within $10^{-4}$ in kinetic
energy, so that column says little beyond the test deck itself. The $\epsilon = 10^{-4}$ column does not separate
the four references. With \texttt{SNAP\_WB\_REF4} on, the $10^{-12}$ ensemble gives 8/10 on the clamp-only reference and
9/10 on the linear one.

Resolution scan with the test deck (nx2 = 4 nx1):
\begin{itemize}
\item the clamp-only reference fails at nx1 128 (34.38 s, in the interior);
\item the linear one fails at nx1 64;
\item this closure reaches 60 s at nx1 32, 64 and 128, with at most 3 consecutive redos.
\end{itemize}

\textbf{The linear form's failure.} The references at $t = 0$ are not odd for any closure. In the straka column the top
cell's $\rho_{\rm ref}$ is within 0.002\% of the cell average with the continuation, against 0.15\% with the clamp,
and the positivity fallback does not fire. The failure develops after about 28 s:
\begin{itemize}
\item on the failing run, the linear $r_{+3}/r_0$ at the top wall falls to 0.35 at 30.9 s and to $-0.15$ at 31.67 s;
\item on the clamp-only run, which survives, the same quantity would reach $-4.9$ at 32.1 s.
\end{itemize}

That the near-zero continuation starves the top cells is a hypothesis; it was not isolated.

\textbf{Not covered.} Non-uniform x1: the continuation is in index, as the binomial is. The MPS tensor path is changed
to the same rule but was not run here, since there is no MPS device.

\end{document}
